\begin{musica}{O Pito da Maria}{Popular Minhota}{Tuna de Engenharia da Universidade do Porto}
\nota{Introdução: C G C G C}

[De todas as capo\ac{G}{ei}ras
Lá da minha fregue\ac{C}{sia},] \nota{x2}
[Não há pito mais bo\ac{G}{ni}to
Que o pito da minha Ma\ac{C}{ria}.] \nota{x2}

\nota{Refrão:}
[\ac{C}{Ó} Maria, dá o pito,
Ó Maria, dá o \ac{G}{cá},
Ó Maria, dá o pito,
Toma lá, dá \ac{C}{cá}.] \nota{x2}

[Que pito tão rechonchudo,
Tão preto e cabeludo,
Havia até quem dissesse
Que o comia com penas e tudo.] \nota{x2}

[Um belo dia porém
Chegou o engenheiro aflito:
“Oh Maria, oh Maria,
Deixaste comer o pito!”] \nota{x2}

[S. Gonçalo de Amarante
Que estás virado p’rá vila,
Virai-vos p’ra outro lado
Que vos dá o Sol na testa.] \nota{x2}

\posicoes{C: x32010, G: 320033}
\end{musica}
